\documentclass[10pt]{article}
\usepackage[utf8]{inputenc}
\usepackage[a4paper,height=20cm,width=13cm]{geometry}
\usepackage{amssymb}
\usepackage{dsfont}
\usepackage{calc}
\usepackage{graphicx}
% \usepackage{pst-node}
\usepackage{fourier}
\usepackage{euscript}
\usepackage{amsmath,amsthm}
\def\phi{\varphi}
\def\theta{\vartheta}
\def\P{\EuScript P}
\def\M{\EuScript M}
\def\B{\EuScript B}
\def\C{\EuScript C}
\def\D{\EuScript D}
\def\U{\EuScript U}
\def\V{\EuScript V}
\def\S{\EuScript S}
\def\sm{\smallsetminus}
\def\niff{\nleftrightarrow}
\def\ZZ{\mathds Z}
\def\NN{\mathds N}
\def\PP{\mathds P}
\def\QQ{\mathds Q}
\def\RR{\mathds R}
\def\<{\langle}
\def\>{\rangle}
\def\E{\exists}
\def\A{\forall}
\def\0{\varnothing}
\def\imp{\rightarrow}
\def\iff{\leftrightarrow}
\def\IMP{\Rightarrow}
\def\IFF{\Leftrightarrow}
\def\range{\textrm{im}}
\def\Mod{\textrm{Mod}}
\def\Aut{\textrm{Aut}}
\def\Th{\textrm{Th}}
\def\acl{\textrm{acl}}
\def\dcl{\textrm{dcl}}
\def\dom{{\rm dom}}
\def\eq{{\rm eq}}
\def\fp{{\rm fp}}
\def\tp{\textrm{tp}}
\def\equivL{\stackrel{\smash{\scalebox{.5}{\rm L}}}{\equiv}}


\def\nonfork{\mathop{\raise0.2ex\hbox{\ooalign{\hidewidth$\vert$\hidewidth\cr\raise-0.9ex\hbox{$\smile$}}}}}

\newcommand{\labella}[1]{{\sf\footnotesize #1}\hfill}
\renewenvironment{itemize}
  {\begin{list}{$\triangleright$}{%
   \setlength{\parskip}{0mm}
   \setlength{\topsep}{0mm}
   \setlength{\rightmargin}{0mm}
   \setlength{\listparindent}{0mm}
   \setlength{\itemindent}{0mm}
   \setlength{\labelwidth}{3ex}
   \setlength{\itemsep}{0mm}
   \setlength{\parsep}{0mm}
   \setlength{\partopsep}{0mm}
   \setlength{\labelsep}{1ex}
   \setlength{\leftmargin}{\labelwidth+\labelsep}
   \let\makelabel\labella}}{%
   \vspace*{-.3\baselineskip}
  \end{list}}
%\def\ssf#1{\textsf{\small #1}}
\newcounter{ex}
\newenvironment{exercise}{\bigskip\addtocounter{ex}{1}\textbf{Esercizio \theex.\ }}{}
\pagestyle{empty}
\parindent0ex
\parskip1ex
\raggedbottom
\def\nsR{{}^*\!\RR}
\def\ssf#1{\textsf{#1}}
\renewcommand{\baselinestretch}{1.3}


\usepackage{fancyhdr}
\pagestyle{fancy}
\lfoot{Teoria dei Modelli a.a.~2026/27}
\rhead{}
\rfoot{{\small\thepage}}
\cfoot{}

%
% INDICARE NOME E COGNOME DI TUTTI GLI AUTORI
%
\lhead{Nome/i Cognome/i}
\begin{document}


\begin{exercise}(Beatrice Calderaro \& Filippo Giugie)
 For any model $M$ and any $b\in\U^z$ prove that the following are equivalent
  \begin{itemize}
    \item [1.] for every $a$, if $a\nonfork_Mb$ then $b\nonfork_Ma$
    \item [2.] for every $\phi(x\,;z)\in L(M)$ there are some $b_1,\dots,b_n\in M^z$ such that \smallskip
    
     \hfil $\phi(M^x;b)\ \ \subseteq\ \ \phi(M^x;b_1)\cup\dots\cup\phi(M^x;b_n).$
  \end{itemize}
\end{exercise}

\begin{exercise}(Martina Apicella \& Andrea Morello)
  The following claim is too strong to be true.
  For every finite coloring of $\NN^{(2)}$ there is an infinite sequence $\bar a=\langle a_i:i<\omega\rangle$ in $\NN$ such that $\fp(\bar a)^{(2)}$ is monochromatic.
  Replace $\fp(\bar a)^{(2)}$ with a (slightly) smaller set that makes the claim true.
\end{exercise}

\begin{exercise}(Mattia Lo Mele \& Daniele Lanciotti)
  Let $M$ be a semigroup $M$ with the property that $a\cdot M$ is infinite for every $a\in M$.
  Prove that for every finite coloring of $M$ 
  there is a sequence $\bar a=\langle a_i:i<\omega\rangle$ with an infinite range such that $\fp\,\bar a$ is monochromatic.
\end{exercise}

\begin{exercise}(Francesco Dam\&	Gabriele Maiolino)
  For every finite coloring of $\NN$ and every $k<\omega$, 
  there is an infinite sequence $\bar a=\<a_i(x):i<\omega\>$ of nonconstant univariate polynomials with coefficients in $\NN$ such that $\big\{a(i): a(x)\in\fp(\bar a),\ i<k\big\}$ is monochromatic.
  Can we require that the $a_i(x)$ have increasing degrees?
\end{exercise}

\end{document}


